\chapter{Citation}
\label{chap:citation}

\section{Citing ActSim}

If you use \actsim in research, internal actuarial analysis, or
educational materials, please cite the package. A general-purpose
citation is:

\begin{quote}
Zhang, J. \emph{ActSim: A Python-Based Actuarial Simulation Package}.
GitHub repository, \url{https://github.com/jzhng105/actsim}.
\end{quote}

A BibTeX-style entry suitable for inclusion in academic work is:

\begin{lstlisting}[style=bashstyle]
@software{ActSim2025,
    title  = {ActSim: A Python package for actuarial risk modeling
              and simulation},
    author = {Juntao Zhang},
    year   = {2025},
    url    = {https://github.com/jzhng105/actsim},
}
\end{lstlisting}

If \actsim is later submitted to a CAS publication, e-Forum paper, or
official research resource, this section will be updated to include
the formal citation.

\section{Citing Dependencies}

\actsim relies on a number of open-source scientific Python libraries.
Users who report results based on \actsim are encouraged to cite the
underlying packages where appropriate, including:

\begin{itemize}[leftmargin=*]
    \item NumPy
    \item SciPy
    \item pandas
    \item matplotlib
    \item seaborn
    \item statsmodels
    \item chainladder (where used for reserving analyses)
\end{itemize}

Citing these dependencies acknowledges the work of the broader
open-source ecosystem on which \actsim depends.
